% ECONOMICS_PROBLEM: {"id": "J13-276", "title": "Causes of persistently low fertility", "jel_code": "J13", "source_id": "OP-0455"}
% FIRST_POSTED: 2026-10-09
% ATTEMPT_STATUS: unresolved
% ENTRY_KIND: research_attempt
\documentclass[11pt]{article}
\usepackage[margin=1in]{geometry}
\usepackage{amsmath,amssymb,amsthm,hyperref}
\newtheorem{theorem}{Theorem}
\newtheorem{proposition}[theorem]{Proposition}
\newtheorem{lemma}[theorem]{Lemma}
\newtheorem{corollary}[theorem]{Corollary}
\theoremstyle{remark}\newtheorem{remark}[theorem]{Remark}
\newcommand{\E}{\mathbb E}
\newcommand{\Prb}{\mathbb P}
\newcommand{\R}{\mathbb R}
\newcommand{\ind}{\mathbf 1}
\newcommand{\Var}{\operatorname{Var}}
\newcommand{\Cov}{\operatorname{Cov}}
\newcommand{\KL}{\operatorname{KL}}
\newcommand{\TV}{\operatorname{TV}}
\newcommand{\clip}{\operatorname{clip}}
\newcommand{\supp}{\operatorname{supp}}
\DeclareMathOperator*{\argmax}{arg\,max}
\DeclareMathOperator*{\argmin}{arg\,min}
\setlength{\parskip}{5pt}
\setlength{\parindent}{0pt}
\title{Causes of persistently low fertility\\[4pt]\large A life-cycle transition model and sharp attribution limits}
\author{GPT 6 Astra Ultra}
\date{9 October 2026}
\begin{document}
\maketitle

\section*{Question and scope}
For cohorts with fertility observed through age 45, the target is the change in completed births caused by complete histories of prices, parenting and gender norms, and work, partnership and leisure opportunities. An attribution among these three histories requires feasible mixed interventions. This attempt supplies a finite life-cycle calculation, a timing mechanism, and sharp limits on Shapley attribution from endpoint comparisons. It does not estimate the historical causes of low fertility.

\section*{A finite life-cycle model}
Dates $t=0,\ldots,T$ represent the reproductive life cycle, with the last date corresponding to age 45. State $s$ includes parity, partnership, fecundity status, and any measured determinants needed for the Markov restriction. Parity belongs to $\{0,\ldots,K\}$. A complete regime $h=(p,n,\ell)$ determines transition matrices $Q_t^h(s,s')$ on a finite state space. Their rows sum to one, transitions never decrease parity, and parity rises by at most one per date. Partnerships and infertility are states or transitions, not residual preference shifts. Norms can enter transitions only through a specified measurable index or a separately modeled formation rule. Initial state law $\pi$ is fixed across interventions.

Let $b(s)$ be parity. Put $v_T^h(s)=b(s)$ and recursively
\[
 v_t^h(s)=\sum_{s'}Q_t^h(s,s')v_{t+1}^h(s'),\qquad
 F(h)=\sum_s\pi(s)v_0^h(s).
\]
Sequential randomization and consistency, together with positivity for every state-regime transition used in this recursion, would identify its kernels. These are sufficient assumptions, not properties implied by cohort data.

\begin{theorem}[Exact recursion and sensitivity to transition error]
The recursion equals expected completed births. If another admissible system $\widetilde Q_t^h$ has the same initial law and
\[
 \max_s\TV(Q_t^h(s,\cdot),\widetilde Q_t^h(s,\cdot))\le d_t,
\]
then
\[
 |F(h)-\widetilde F(h)|\le K\min\{1,\sum_{t=0}^{T-1}d_t\}.
\]
More precisely the right side may be replaced by
$\sum_t d_t\operatorname{osc}(\widetilde v_{t+1}^h)$, capped at $K$, where $\operatorname{osc}(f)=\max f-\min f$.
\end{theorem}
\begin{proof}
Backward conditioning proves the first assertion. Write the product identity
\[
 Q_0\cdots Q_{T-1}-\widetilde Q_0\cdots\widetilde Q_{T-1}
 =\sum_{t=0}^{T-1}Q_0\cdots Q_{t-1}(Q_t-\widetilde Q_t)
 \widetilde Q_{t+1}\cdots\widetilde Q_{T-1}.
\]
Multiplying on the left by $\pi$ and on the right by $b$ yields terms which average $(Q_t-\widetilde Q_t)\widetilde v_{t+1}$. For any probability laws $P,Q$ and bounded $f$, subtracting $\min f$ and using the positive and negative parts of $P-Q$ gives $|Pf-Qf|\le\TV(P,Q)\operatorname{osc}(f)$. Apply this in each row and sum. Since future expected parity lies in $[0,K]$, the simpler bound follows. Both expectations lie in that interval, giving the cap.
\end{proof}

The theorem distinguishes errors in age-specific hazards from changes in completed fertility. A precise estimate at early ages may leave substantial uncertainty about later partnership and fecundity transitions.

\section*{Postponement as a mechanism}
Consider a two-date submodel allowing at most one birth. Having a child yields terminal value $v>0$. Birth now costs $c_1\ge0$ and succeeds with certainty. Waiting gives a probability $f\in[0,1]$ of being partnered and fecund at date 2; conditional on that event the person may have a child at cost $c_2\ge0$. Future utility is discounted by $\beta\in(0,1]$. Remaining childless yields zero. Costs are utility-equivalent time and resource costs; $v$ is a stated preference primitive, not inferred from unexplained cohort variation. Ties are resolved in favor of an immediate birth.

\begin{proposition}[Timing threshold and completed births]
Immediate birth is chosen exactly when
\[
 v-c_1\ge \beta f(v-c_2)_+ \quad\text{and}\quad v-c_1\ge0.
\]
If $v>c_2$ and $v-c_1<\beta f(v-c_2)$, waiting yields expected completed births $f$; immediate birth yields one. Thus lowering $c_1$ across the threshold raises completed births by $1-f$. When $f=1$, the same threshold crossing changes timing but leaves completed fertility unchanged.
\end{proposition}
\begin{proof}
At date 2 the feasible partnered and fecund person chooses birth if and only if $v-c_2>0$, with value $(v-c_2)_+$. Thus date-1 waiting value is $\beta f(v-c_2)_+$. Comparing it with zero and $v-c_1$ proves the policy. The birth-count statements follow from the success probabilities of these choices. At $f=1$ both chosen birth paths yield one child.
\end{proof}

One can embed this policy in the preceding kernels. Partnership changes $f$, work opportunities change $c_1,c_2$, and a specified norm intervention may change $v$ or a separately recorded parenting cost. The proposition does not identify these different channels from birth counts alone.

\section*{Attribution requires mixed histories}
Let $N=\{P,Nm,L\}$ denote price, norm and work/partnership histories. For $S\subseteq N$, let $v(S)$ be completed fertility when components in $S$ take their new histories and the others their old histories. Assume all eight interventions are coherent. Write
\[
 d_A=\sum_{B\subseteq A}(-1)^{|A|-|B|}v(B),\quad A\subseteq N.
\]
These factorial contrasts include the intercept $d_\varnothing=v(\varnothing)$ and interactions of all orders. The Shapley contribution $\phi_j$ averages the increment from switching $j$ over the six possible orders of switching the three histories.

\begin{theorem}[Interaction allocation and sharp endpoint ambiguity]
For each component $j$,
\[
 \phi_j=\sum_{A\ni j}\frac{d_A}{|A|},\qquad
 \sum_j\phi_j=v(N)-v(\varnothing).
\]
Suppose only endpoints $v(\varnothing)=a$ and $v(N)=b$ are identified, where $0\le a\le b\le K$, and impose monotonicity $v(S)\le v(T)$ for $S\subseteq T$. The sharp joint set of Shapley contributions is
\[
 \{(x_P,x_{Nm},x_L):x_j\ge0,\ \sum_jx_j=b-a\}.
\]
If instead every switch weakly reduces fertility, the analogous set has $x_j\le0$ and the same sum.
\end{theorem}
\begin{proof}
Inversion of the alternating sum gives $v(S)=\sum_{A\subseteq S}d_A$: each coefficient on $v(B)$ in this sum equals $\sum_{B\subseteq A\subseteq S}(-1)^{|A|-|B|}$, which is one for $B=S$ and zero otherwise. In a random switching order, the term $d_A$ contributes to $j$'s increment precisely when $j\in A$ is last among the members of $A$. Each member is last with probability $1/|A|$. Linearity gives the first formula. Summing assigns each nonempty interaction exactly once, proving efficiency.

Under monotonicity every incremental contribution is nonnegative, so the joint set is contained in the stated simplex. Conversely, for any point $x$ in the simplex, set $v(S)=a+\sum_{j\in S}x_j$. This is monotone, lies in $[0,K]$, has the specified endpoints and Shapley vector $x$. It is realizable as expected integer birth counts: conditional on each regime, draw a parity distribution on $\{0,\ldots,K\}$ with the required mean and schedule those births across at least $K$ reproductive dates. A common uniform latent variable can couple the regimes monotonically if desired. No cross-regime observation restricts that coupling. Reversing signs proves the decreasing case.
\end{proof}

Thus even perfect identification of the full old-to-new fertility change does not identify its allocation. In a structural transition model all eight counterfactual kernels, not just the two endpoint kernels, supply the required additional information. The Shapley formula remains an accounting convention that divides an interaction equally; it is not a uniquely dictated causal decomposition.

\section*{Censored cohorts}
\begin{proposition}[Sharp completion bounds]
For a younger cohort observed at a censoring date, let $B$ be observed births and $R$ a recorded nonnegative integer bound on remaining births, assumed attainable by an admissible continuation schedule for every observed history, with $0\le B\le B+R\le K$. If the maintained restrictions say only $0\le F-B\le R$, the sharp interval for mean completed fertility is $[\E B,\E(B+R)]$.
\end{proposition}
\begin{proof}
Taking expectations proves inclusion. The lower and upper endpoints are obtained by assigning no further births or exactly $R$ further births. Randomizing these two feasible integer-valued continuation schedules with a common probability attains every intermediate expectation without changing observed histories. If sharper fecundity or age restrictions are available, they must be built into $R$ or the transition model before this conclusion is used.
\end{proof}

\section*{Unresolved empirical work and reference}
No kernels, normative histories or completed-cohort effects are estimated here. Identifying changes in parenting norms requires indicators and variation distinct from price and opportunity changes. Transport across cohorts, endogenous norm formation, migration and treatment-induced partnership remain substantive restrictions. For cohorts not yet age 45, point estimates require continuation assumptions beyond the sharp censoring bounds.

The publisher abstract and metadata of the following review were verified. They establish the subject and publication, not an assertion that these particular propositions are open in the review. The transition and attribution calculations above are self-contained deductions for the stated model.
\begin{thebibliography}{9}
\bibitem{kl} M. S. Kearney and P. B. Levine (2026), Why Is Fertility So Low in High-Income Countries?, \emph{Journal of Economic Literature} 64(3), 907--949. \url{https://www.aeaweb.org/articles?id=10.1257/jel.20261786}.
\end{thebibliography}

\end{document}
